\scanpage{15}
$n=2)$ \qquad $U=X_f$, $\left\{\begin{smallmatrix}f\in\unsure{\operatorname{r}A|\mathfrak r A|\mathfrak m_A},\\ f\ne0\end{smallmatrix}\right.$ \qquad $F=\underline O_{X,U}$. Soit $X'=X-x$, $Y=V(f)$, \inserted{$Y'=Y\cap X'$} \margin{$\deleted{\unsure{F=|Y=}X'}$}

$x_i$ les pts gén. des composantes de $Y$, \unsure{\ldots}.

\[
\begin{aligned}
H_x^0(F),H_x^1(F)&=\text{Ker et Coker de }\underset{\shortparallel\,0}{H^0(X,F)}
 \longrightarrow\underset{\shortparallel\,0}{H^0(X',F)},\quad\text{nuls}.\\
H_x^i(F)&=0\quad\text{pour }i>2\quad\text{pour raisons de dim ch.}
\end{aligned}
\]

Reste : calculer $H_x^2(F)$ et comparer son dual à $A_f=\unsure{H^0(A,\underline O_U)|H^0(U,\underline O_U)}$.

\unsure{Or.|Car.|Comme} \deleted{\unsure{La suite exacte}}\deleted{\unsure{ici on a}} \unsure{puisque|pour} $H^i(X,F)=0$ pour \deleted{$\unsure{\forall|i|H}$} $i\ne0$, \uuline{$\unsure{\forall|\#} F$},
\[
H_x^2(F)\simeq H^1(X',F)\,.
\]

Soit \deleted{$\unsure{G=i_{\ast}F|G=\underline O_Y|\underline O_Y\,F}$} $i:Y'\longrightarrow X'$ l'injection, et $G=i^{-1}\underline O_X$, on a une suite exacte
\[
0\longrightarrow F\longrightarrow\underline O_{X'}\longrightarrow G\longrightarrow0\,,
\]
$G$ concentré sur l'ensemble \uline{fini discret} $Y'$. D'où par suite exacte de cohomologie

\diagram{\begin{tikzcd}[ampersand replacement=\&,column sep=small,row sep=tiny]
0\arrow[r]\& H^0(X',\underline O_{X'})\arrow[r,"\alpha^0"]\& H^0(Y',G)\arrow[r,"\partial"]\& H^1(X',F)\arrow[r,"\alpha^1"]\& H^1(X',\underline O_{X'})\arrow[r]\&0\\
\&\shortparallel\&\shortparallel\&\shortparallel\&\shortparallel\\
\& A\&\displaystyle\prod_i\underline O_{x_i}\& H_x^2(F)\&I\\
\&\&\shortparallel\\
\&\&\displaystyle\underset{g}{\varinjlim}\,Ag^{-1}
\end{tikzcd}}

\margin{$g$ pris \unsure{mt|mt.|nul} entre eux, i.e. \unsure{étr. que|étr. à|étranger à} : $f$, i.e. $(f,g)$ engendre un idéal primaire pour \uline{$\mathfrak m$}}

Prenant les duels à valeurs dans $I$, on trouve

\diagram{\begin{tikzcd}[ampersand replacement=\&,column sep=small,row sep=tiny]
0\& I\arrow[l]\& \displaystyle\underset{g}{\varprojlim}\,{}_g I\arrow[l,"\alpha^{0\ast}"']\& D(H_x^2(F))\arrow[l,"\partial^\ast"']\& A\arrow[l,"\alpha^{1\ast}"']\arrow[dl,dashed,"\text{en}\,A\,\text{\unsure{complet|en lui|modules}}" description]\&0\arrow[l]\\
\&\shortparallel\&\&A_f\arrow[u,"\varphi"]\arrow[ul,dashed,"{\unsure{d|d_g}}"']\\
\& {}_1 I\&\&\shortparallel\\
\&\&\&\displaystyle\underset{m}{\varinjlim}\,Af^{-m}
\end{tikzcd}}

l'application $A_f\longrightarrow\underset{g}{\varprojlim}\,{}_g I$ s'explicite par
\[
d_g:\quad h\leadsto [h]\left[\frac{1}{g}\right]
\]
le noyau de $d_g$ est alors l'image de $A\longrightarrow A_f$, donc $\varphi$ \uline{est injective}. \uline{Cependant l'image de} $d$ \uline{est} $\ne\operatorname{Ker}\alpha^{0\ast}$, $\varphi$ \uline{n'est pas surjective} ; la raison en est qu'il faut prendre des \uline{\unsure{duels|duals|duaux} topologiques} et non seulement
% archive sheet number: 14

